\begin{problem}[Local fractions and the basis-sheaf construction agree]\label{prob:vi17-p19}
Compare two constructions of $\OO_{\Spec A}$: (i) locally representable families in the stalks $A_{\mfp}$, and (ii) the sheaf determined on the basis $D(f)$ by $A_f$.  Prove they are canonically isomorphic.
\end{problem}

\ifdefined\FullProblemDossiers
\paragraph{Why this problem matters.}
The two standard definitions emphasize different ideas---germs versus localization on a basis---and should produce the same object.

\paragraph{Useful ideas.}
Use the principal-open theorem for construction (i), then apply basis reconstruction.

\paragraph{Guided hints.}
\begin{enumerate}
\item Construction (i) has $\OO(D(f))\cong A_f$.
\item The restriction maps coincide with localization maps.
\item Invoke uniqueness of a sheaf reconstructed from basis data.
\end{enumerate}

\begin{solution}
For the locally representable-family construction, Theorem~\ref{thm:vi17-basic-open} provides canonical identifications $\OO(D(f))\cong A_f$.  Under these identifications, restriction to $D(fg)\subseteq D(f)$ is the localization map $A_f\to A_{fg}$.  These are exactly the data used in the basis construction.  Since principal opens form a basis stable under finite intersection, the basis-reconstruction theorem gives a unique isomorphism of sheaves extending the identity on every $A_f$.
\end{solution}

\paragraph{What happened mathematically.}
The structure sheaf can be viewed either as a sheaf of locally varying germs or as the sheaf generated by the localization presheaf on principal opens.

\paragraph{Extensions.}
Carry out the same comparison for the sheaf $\widetilde M$ associated to an $A$-module.

\paragraph{Related problems.}
VI.17.P18 and VI.17.P20 develop neighboring aspects of the same localization-sheaf calculus.

\paragraph{Provenance.}
\textbf{New synthesis of the presheaf/basis material from VI/12 with the current chapter.}
\else
\paragraph{Provenance.} \textbf{New synthesis of the presheaf/basis material from VI/12 with the current chapter.}
\fi
